\chapter{\texttt{tricycle\_env.py}: the environment}
\label{ch:env}

\textbf{Role:} the behaviour. It fills in the six hooks that
\code{DirectRLEnv} leaves abstract (\code{_setup_scene},
\code{_pre_physics_step}, \code{_apply_action}, \code{_get_observations},
\code{_get_rewards}, \code{_get_dones}) plus the reset hook
\code{_reset_idx}. Every method operates on tensors with a leading env
dimension. 120 lines. Read Section~\ref{sec:step_detail} first: the order in
which \code{DirectRLEnv.step} calls these methods is what makes them correct.

\section{Header, imports, constructor}

\codepart{\srcroot/luna_hifi_tasks/tricycle/tricycle_env.py}{1}{31}

\begin{explain}
\lis{1}{5} What the file exercises. ``Reward is per-step forward
  displacement, which sums to total distance'' is the key design fact: the
  episode return equals metres driven.
\li{7} Modern annotation syntax.
\li{9} PyTorch; every quantity is a torch tensor on the GPU.
\li{11} \code{Articulation}, used only as a type annotation on line 38.
\li{12} The base class.
\li{14} \code{MPMObject}, used as a type annotation on line 39.
\li{16} Joint names from the model file.
\li{17} The cfg class, for the \code{cfg:} annotation and for the
  constructor signature.
\li{20} The env class. \code{DirectRLEnv} is also a \code{gym.Env}, so
  \code{gym.make} can build it.
\li{21} A class-level annotation telling type checkers (and readers) that
  \code{self.cfg} is a \code{TricycleEnvCfg}, so \code{self.cfg.max_steer}
  is known to exist.
\li{23} The constructor signature Isaac Lab expects: the cfg, an optional
  render mode, and pass-through keyword arguments.
\li{24} Runs everything in Section~\ref{sec:step_env}: simulation context,
  scene, our \code{_setup_scene}, physics start, buffers. When it returns,
  \code{self.car} exists and its data is live.
\li{26} Look up the steering joint's index by name.
  \code{find_joints} returns \code{(indices, names)}; we keep the indices.
  The result is a list with one element.
\li{27} Likewise the two rear-wheel indices, in the order given
  (\code{rear_left}, \code{rear_right}).
\li{28} Fail loudly at start-up if the asset does not contain exactly those
  joints, rather than producing wrong tensors later. This assertion is what
  would fire if the USD were regenerated with renamed joints.
\li{30} A buffer for the last action, shape $(N, 2)$, on the simulation
  device. It is reused every step to avoid allocations and is also part of
  the observation.
\li{31} A buffer for each car's $x$ position at the previous control step,
  shape $(N,)$, used by the reward.
\end{explain}

\begin{isaacbox}{\code{self.num_envs} and \code{self.device}}
Both come from \code{DirectRLEnv}: \code{num_envs} is
\code{cfg.scene.num_envs} after command-line overrides, and \code{device}
is \code{cfg.sim.device} (\code{cuda:0}). All buffers you create must live
on \code{self.device} or every step pays a host--device copy.
\end{isaacbox}

\section{Scene hook}

\codepart{\srcroot/luna_hifi_tasks/tricycle/tricycle_env.py}{33}{39}

\begin{explain}
\li{35} Called by \srcpath{DirectRLEnv.__init__} \emph{after}
  \srcpath{InteractiveScene(self.cfg.scene)} has spawned everything and
  \emph{before} the simulation starts.
\lis{36}{37} The trap this file used to fall into: constructing an
  \code{Articulation(cfg)} here would try to spawn the prim a second time and
  fail with ``A prim already exists at path''. Isaac Lab's tutorial tasks
  create assets here because their scene cfgs do not list them; ours does.
\li{38} Fetch the car by its scene entry name (line 181 of the cfg). The
  annotation makes later attribute access readable in an editor.
\li{39} Fetch the soil. It is only needed for \code{reset}.
\end{explain}

\section{Actions}

\codepart{\srcroot/luna_hifi_tasks/tricycle/tricycle_env.py}{41}{50}

\begin{explain}
\li{43} Called once per control step with the raw policy output, shape
  $(N, 2)$, before any physics.
\li{44} Clamp to $[-1,1]$ and store \emph{in place} into the buffer
  (\code{[:]}). The policy is a Gaussian and can emit values outside the
  range; clamping keeps the scales on the next lines meaningful. In-place
  assignment keeps the same tensor object that
  \code{_get_observations} concatenates.
\li{46} Called once per \emph{physics} sub-step (twice per control step with
  \code{decimation=2}), or once per control step when Newton folds the
  decimation loop into one call. Either way it writes the same targets.
\li{47} Throttle in rad/s: column 0 (kept as an $(N,1)$ slice with
  \code{0:1}) times \code{max_wheel_speed}.
\li{48} Steering in rad: column 1 times \code{max_steer}.
\li{49} Write the same speed to both rear wheels.
  \code{expand(-1, 2)} turns $(N,1)$ into a view of shape $(N,2)$ without
  copying; \code{joint_ids} selects the two rear joints found in the
  constructor. Because the rear actuators have zero stiffness, this target
  alone defines the drive torque.
\li{50} Write the steering angle as a position target to the steering
  joint. Targets are buffered by the \code{Articulation} and flushed to the
  solver by \srcpath{scene.write_data_to_sim()} right after this method
  (Figure~\ref{fig:step}).
\end{explain}

\begin{whybox}{no differential steering}
Both rear wheels get the same speed target and only the front wheel steers.
The car therefore cannot ``skid-steer'' by spinning the rear wheels at
different speeds; every turn comes from the fork. It keeps the action space
at two numbers and makes the steering observation meaningful.
\end{whybox}

\section{Observations}

\codepart{\srcroot/luna_hifi_tasks/tricycle/tricycle_env.py}{52}{68}

\begin{explain}
\li{54} Called at the end of every step, after resets. Must return a dict
  with at least a \code{"policy"} entry whose tensor matches
  \code{observation_space}.
\li{55} A short alias for the car's data container. Its fields are
  recomputed by \code{scene.update} each sub-step.
\li{56} Concatenate along the last dimension \ldots
\li{58} \ldots\ the root linear velocity in the body frame, $(N,3)$: forward
  speed, sideways speed, vertical speed as the car feels them;
\li{59} the root angular velocity in the body frame, $(N,3)$: roll, pitch,
  yaw rates;
\li{60} gravity projected into the body frame, $(N,3)$: a compact encoding
  of roll and pitch that does not need a heading (yaw is deliberately not
  observed, so the policy is heading-independent);
\li{61} the steering angle, $(N,1)$;
\li{62} the steering rate, $(N,1)$;
\li{63} the two rear-wheel speeds, $(N,2)$, which together with forward
  speed let the policy infer slip;
\li{64} the previous action, $(N,2)$, which helps the policy produce smooth
  commands;
\lis{65}{67} \ldots\ for a total of $3+3+3+1+1+2+2 = 15$, matching the cfg.
\li{68} Return under the \code{"policy"} key. The \code{rsl_rl} wrapper
  feeds the same group to actor and critic because no \code{obs_groups} are
  configured.
\end{explain}

\begin{whybox}{no position in the observation}
The car's $x$ position is not observed. The reward is displacement, not
position, so the optimal action does not depend on where along the strip the
car is; leaving position out removes a feature the policy could only misuse
and keeps the observation translation-invariant.
\end{whybox}

\section{Reward}

\codepart{\srcroot/luna_hifi_tasks/tricycle/tricycle_env.py}{70}{76}

\begin{explain}
\li{72} Called once per step, after \code{_get_dones} and before resets.
  Returns $(N,)$.
\li{73} Each car's world $x$ minus its env origin's $x$: the env-local
  forward position, so cars in different envs are comparable.
\li{74} The displacement since last step, with any NaN or infinity replaced
  by zero. If the physics diverged this step, the reward is zero rather than
  garbage, and \code{_get_dones} will terminate the env.
\li{75} Remember the position for next step (in place).
\li{76} Clamp to $\pm0.1$~m per 20~ms step, i.e.\ $\pm5$~m/s. A car
  driving at its top speed of about 2~m/s earns 0.04 per step; the clamp
  only bites on non-physical jumps, so a diverging simulation cannot pay
  out a giant reward before it is terminated. Summed over an episode this
  is metres travelled forward.
\end{explain}

\section{Terminations}

\codepart{\srcroot/luna_hifi_tasks/tricycle/tricycle_env.py}{78}{93}

\begin{explain}
\li{80} Returns two boolean $(N,)$ tensors: \code{terminated} and
  \code{time_out} (Section~\ref{sec:term_vs_timeout}).
\li{81} Time-out when the per-env step counter reaches
  \code{max_episode_length - 1}, i.e.\ after 100 control steps. This is the
  idiom every Isaac Lab direct task uses (the counter was already
  incremented for this step).
\li{83} Env-local position of each car, $(N,3)$.
\li{84} Body-frame linear velocity, $(N,3)$.
\li{86} ``Invalid'' if any component of position or velocity is NaN or
  infinite: the physics has blown up.
\li{87} \ldots\ or the car is more than 50~m from its origin;
\li{88} \ldots\ or moving faster than 50~m/s. Both are impossible for a real
  car and catch a diverging solver a step or two before it produces NaNs.
\li{90} Off the strip sideways: $|y|$ beyond \code{max_lateral_drift}.
\li{91} Flipped: body-frame gravity $z$ below the threshold (orange box in
  Chapter~\ref{ch:env_cfg}).
\li{93} Terminate on any of the three; time out separately. An env can be
  both terminated and timed out in the same step; \code{rsl_rl} treats
  \code{terminated} as taking priority.
\end{explain}

\begin{whybox}{why the invalid-state termination exists}
The runbook records a run where the reward exploded to $10^8$: a coupled
step had diverged, the car was flung away at absurd speed, and the
displacement reward faithfully paid for it. Terminating on non-finite or
absurd state, together with the reward clamp, turns a diverging env into a
short, low-reward episode instead of a poisoned training batch.
\end{whybox}

\section{Reset}

\codepart{\srcroot/luna_hifi_tasks/tricycle/tricycle_env.py}{95}{120}

\begin{explain}
\li{97} Called with the indices of the envs that ended this step, and once
  with all indices at start-up (from \code{DirectRLEnv.reset}).
\lis{98}{99} Defensive: if ever called with \code{None}, reset everything.
  \code{DirectRLEnv} never passes \code{None}, so this is a safety net (the
  previous version referenced an attribute that does not exist on
  \code{Articulation}; it has been replaced with \code{torch.arange}).
\li{100} The base implementation: \srcpath{scene.reset(env_ids)} (which calls
  \code{reset} on every asset and clears their internal buffers), applies
  any reset-mode events, and zeroes \srcpath{episode_length_buf[env_ids]}.
\li{102} Start from the default root state of the selected envs (position,
  quaternion, linear and angular velocity: 13 numbers per env, from the
  cfg's \code{init_state}). \code{clone()} because the next line modifies
  it.
\li{103} Add each env's origin to the position, turning env-local spawn
  coordinates into world coordinates.
\li{104} Write pose (7 numbers: position and quaternion) to the simulator
  for those envs.
\li{105} Write velocity (6 numbers) for those envs: all zero.
\lis{106}{111} Write joint positions and velocities back to their defaults
  (all zero) for all joints (\code{None}) of those envs.
\lis{113}{116} What \code{soil.reset} does and does not do: it rewrites the
  default positions and zero velocities of every particle belonging to those
  envs (Isaac Lab's \code{MPMObject.reset} calls
  \srcpath{write_nodal_state_to_sim_index} with the default nodal state). It
  does not clear the solver's internal history; Isaac Lab offers
  \srcpath{NewtonMPMManager.reset_solver_state} for an exact reset, and the
  tricycle has not needed it.
\li{117} Reset the soil of those envs.
\li{119} Zero the stored actions for those envs, so the observation's
  ``previous action'' is zero at episode start.
\li{120} Initialise the reward's memory to the env-local spawn $x$ (which is
  0), so the first step's displacement is measured from the spawn point and
  not from wherever the previous episode ended.
\end{explain}

\begin{isaacbox}{index-based versus mask-based writes}
Isaac Lab develop offers two flavours of every asset write:
\code{write_root_pose_to_sim(data, env_ids)} (index-based, used here) and
\code{write_root_pose_to_sim_mask(data, env_mask)} (mask-based, avoids a
GPU--CPU sync when many envs reset). For a handful of envs the index form is
simplest and its cost is negligible.
\end{isaacbox}
